% Generated draft: requires mathematical and source review.
\documentclass[11pt]{article}
\usepackage[margin=25mm]{geometry}
\usepackage{fontspec}
\setmainfont{Latin Modern Roman}
\usepackage{amsmath,amssymb,mathtools,unicode-math}
\setmathfont{Latin Modern Math}
\usepackage{enumitem,needspace,xurl,hyperref}
\setlength{\parindent}{0pt}
\setlength{\parskip}{5pt}
\setlength{\emergencystretch}{3em}
\newcommand{\N}{\mathbb{N}}
\newcommand{\Z}{\mathbb{Z}}
\newcommand{\Q}{\mathbb{Q}}
\newcommand{\R}{\mathbb{R}}
\newcommand{\C}{\mathbb{C}}
\newcommand{\F}{\mathbb{F}}
\newcommand{\A}{\mathbb{A}}
\newcommand{\PP}{\mathbb{P}}
\newcommand{\E}{\mathbb{E}}
\newcommand{\eps}{\varepsilon}
\newcommand{\dd}{\,\mathrm{d}}
\begin{document}
\section*{UnsolvedMath ID: 3100007}
\textbf{Source problem code:} catalogue entry
\subsection{Definitions and mathematical statement}
Let $V\subset\R^3$ consist of $n$ distinct points, and let $L$ be a finite set of affine lines. A triple $a,b,c\in V$ spans the convex triangle $\operatorname{conv}\{a,b,c\}$ when the points are noncollinear. The source assumes that every triangle spanned by three points is pierced by some line in $L$ and asks for a linear lower bound on $|L|$.

The intended asymptotic assertion is that there exist absolute constants $c>0$ and $n_0$ such that every admissible configuration with $n\geq n_0$ satisfies
\[
 \bigl[\forall\{a,b,c\}\subseteq V\text{ forming a triangle},\ 
 \exists\ell\in L:\ell\text{ pierces that triangle}\bigr]
 \quad\Longrightarrow\quad |L|\geq cn.
\]
Here admissible configurations and the meaning of pierces must be recovered from the source. The record does not say whether piercing means meeting the relative interior transversely, meeting the closed triangle, or another incidence condition, and it imposes no general-position hypothesis.

These omissions are essential. If only nondegenerate triangles are quantified, an entirely collinear point set has no such triangles and the hypothesis is vacuous. If intersections on boundaries or line containment are allowed, degenerate triples and shared incidences change the problem again. Thus the literal arbitrary-point formulation does not determine the intended positive lower-bound conjecture. No general-position or strict-piercing convention is silently inserted.

The source reports a lower bound of order $n^{1/2}$ as the best known in its intended setting. The assessment likewise calls the linear bound conjectural and says it did not verify an improvement after approximately 2020. The applicable incidence/general-position hypotheses for that reported bound are not supplied.
\subsection{Short English statement}
Solymosi's question seeks a linear number of lines whenever every point-spanned triangle in three-dimensional space must be pierced. The source claims only a square-root lower bound was known, but omits the degeneracy and piercing conventions needed for either bound. The October 2020 snapshot's open label and later unsuccessful verification are dated source evidence, not an independent current-status conclusion.
\subsection{Sources}
\begin{itemize}
\item UnsolvedMath permanent ID 3100007, source code catalogue entry, Solymosi's line-piercing triangle problem. Registry curation and metadata: Ulam AI and UnsolvedMath Contributors, CC BY 4.0; updated 2026-08-06.
\item Cooper, \emph{Combinatorial Problems I Like}, Geometry, source-order bullet 7, October 2020 snapshot: \url{https://people.math.sc.edu/cooper/combprob.html}, supplied HTML extraction accessed 2026-07-29. The author-maintained snapshot labels it open, but later status was not checked item by item. Redistribution permission was not established, as recorded in the supplied rights-review note.
\item Solymosi and unspecified related triangle-piercing work are mentioned by the supplied assessment without complete papers or theorem statements. No missing citation is inferred.
\end{itemize}
\end{document}
